














\documentclass[12pt,a4paper,oneside]{article}

\usepackage[T1]{fontenc}
\usepackage[utf8]{inputenc}
\usepackage[brazilian]{babel}
\usepackage{newtxtext}
\usepackage{newtxmath}
\usepackage[section]{placeins}
\usepackage[a4paper,top=2.5cm,bottom=2.5cm,left=3cm,right=3cm]{geometry}
\usepackage{setspace}
\usepackage{indentfirst}
\usepackage{titlesec}
\usepackage{caption}
\usepackage{booktabs}
\usepackage{array}
\usepackage{multirow}
\usepackage{graphicx}
\usepackage[colorlinks=true,
            urlcolor=blue!60!black,
            linkcolor=black,
            citecolor=black,
            breaklinks=true]{hyperref}
\usepackage{url}
\usepackage{xurl}
\Urlmuskip=0mu plus 1mu\relax
\usepackage{float}
\usepackage{tcolorbox}

\graphicspath{{figures-pbf/}}

\onehalfspacing
\setlength{\parindent}{1.25cm}
\setlength{\parskip}{0pt}

\titleformat{\section}{\normalfont\bfseries\large\uppercase}{\thesection}{0.6em}{}
\titleformat{\subsection}{\normalfont\bfseries\normalsize}{\thesubsection}{0.6em}{}
\titlespacing*{\section}{0pt}{18pt}{8pt}
\titlespacing*{\subsection}{0pt}{12pt}{6pt}

\captionsetup[table]{labelfont=bf,name=Tabela,labelsep=endash,
                     position=above,justification=centering,font=small}
\captionsetup[figure]{labelfont=bf,name=Figura,labelsep=endash,
                      position=below,justification=centering,font=small}
\captionsetup{skip=6pt}

\newcommand{\rs}{R\$\,}

% Stamp de compilação injetado pelo CI (deploy-pages.yml).
% Para builds locais, este arquivo default é usado — basta um pdflatex direto
% e a capa exibe "(dev local)". Em CI, este arquivo é sobrescrito antes da
% compilação com timestamp BRT + SHA do commit que disparou o deploy.
\newcommand{\COMPILESTAMP}{compila\c{c}\~ao local}
\newcommand{\COMPILESHA}{dev}


\begin{document}


\begin{center}{\LARGE\bfseries Três regimes, um programa: documentação, identificação causal e equidade do Bolsa Família\par}\end{center}

\begin{center}{\large\itshape Three regimes, one program: documentation, causal identification, and equity of Bolsa Família\par}\end{center}

\begin{center}{\large\itshape Tres regímenes, un programa: documentación, identificación causal y equidad de Bolsa Família\par}\end{center}













\vspace{6pt}



\begin{center}{\footnotesize\itshape [Identificação dos autores removida para avaliação por pares duplo-cega.]\par}\end{center}

\section*{Resumo}
\begin{singlespace}\small\noindent
Este trabalho analisa o Bolsa Família entre 2013 e 2025, integrando três dimensões usualmente tratadas em separado: documentação reproduzível dos microdados, identificação causal das transições institucionais e equidade distributiva. Consolidam-se microdados mensais do Portal da Transparência (mais de 2,2 bilhões de registros) em painel estado$\times$ano, via arquitetura \textit{medallion}, com deflação pelo IPCA (base dezembro de 2021). A identificação causal articula diferenças-em-diferenças e efeitos fixos de dois sentidos sobre a Medida Provisória 1.061/2021 e a Lei 14.601/2023, com inferência robusta por \textit{wild-cluster bootstrap}, testes de tendências paralelas, placebo e \textit{leave-one-out}. A equidade é medida pelo índice de Kakwani, ordenado pelo IDH municipal, e por um índice de necessidade. Um \textit{benchmark} internacional compara o programa a quatro transferências condicionadas latino-americanas em dólares PPC. Em termos reais, o gasto cresceu 283\,\% entre 2018 e 2024 e o valor por beneficiário, 181\,\%. O efeito causal é grande sob erros HC3, mas indistinguível de zero sob \textit{bootstrap} robusto, e o teste de tendências paralelas é rejeitado --- ilustrando o ônus identificacional com poucos \textit{clusters}. O índice de Kakwani confirma a progressividade do programa, porém decrescente entre regimes; o índice de necessidade revela que estados líderes em valor per capita estão subatendidos frente à intensidade local da pobreza. O artigo publica o \textit{pipeline}, os dados e os testes automatizados, viabilizando reprodução integral.

\vspace{6pt}\noindent\textbf{Palavras-chave:} Bolsa Família; transferência condicionada de renda; avaliação de políticas públicas; equidade distributiva; reprodutibilidade.
\end{singlespace}

\vspace{6pt}
\section*{Abstract}
\begin{singlespace}\small\noindent
This paper analyzes Brazil's Bolsa Família program between 2013 and 2025, integrating three dimensions usually treated separately: reproducible microdata documentation, causal identification of the institutional transitions, and distributive equity. Monthly microdata from the federal Transparency Portal (over 2.2 billion records) are consolidated into a state-by-year panel through a \textit{medallion} architecture, deflated by the IPCA index (December 2021 base). Causal identification combines difference-in-differences and two-way fixed-effects estimators on Provisional Measure 1.061/2021 and Law 14.601/2023, with robust inference via \textit{wild-cluster bootstrap}, parallel-trends tests, placebo, and \textit{leave-one-out} checks. Equity is measured with the Kakwani index ranked by the municipal Human Development Index and with a need index. An international benchmark compares the program with four Latin American conditional cash transfers in PPP dollars. In real terms, spending grew 283\,\% between 2018 and 2024 and the per-beneficiary value grew 181\,\%. The causal effect is large under HC3 errors but indistinguishable from zero under robust \textit{bootstrap}, and the parallel-trends test is rejected --- illustrating the identification burden with few \textit{clusters}. The Kakwani index confirms the program's progressivity, although it declines across regimes; the need index reveals that states leading in per-capita value are underprovisioned relative to local poverty intensity. The paper releases the pipeline, dataset, and automated tests, enabling full reproduction.

\vspace{6pt}\noindent\textbf{Keywords:} Bolsa Família; conditional cash transfers; public policy evaluation; distributive equity; reproducibility.
\end{singlespace}

\vspace{6pt}
\section*{Resumen}
\begin{singlespace}\small\noindent
Este trabajo analiza el programa Bolsa Família entre 2013 y 2025 integrando tres dimensiones habitualmente tratadas por separado: documentación reproducible de los microdatos, identificación causal de las transiciones institucionales y equidad distributiva. Se consolidan microdatos mensuales del Portal de la Transparencia (más de 2,2 mil millones de registros) en un panel estado$\times$año, mediante una arquitectura \textit{medallion}, con deflación por el IPCA (base diciembre de 2021). La identificación causal combina diferencias-en-diferencias y efectos fijos de dos vías sobre la Medida Provisional 1.061/2021 y la Ley 14.601/2023, con inferencia robusta por \textit{wild-cluster bootstrap}, pruebas de tendencias paralelas, placebo y \textit{leave-one-out}. La equidad se mide con el índice de Kakwani, ordenado por el IDH municipal, y con un índice de necesidad. Un \textit{benchmark} internacional compara el programa con cuatro transferencias monetarias condicionadas latinoamericanas en dólares PPA. En términos reales, el gasto creció 283\,\% entre 2018 y 2024 y el valor por beneficiario, 181\,\%. El efecto causal es grande con errores HC3, pero indistinguible de cero bajo \textit{bootstrap} robusto, y la prueba de tendencias paralelas se rechaza, ilustrando la carga de identificación con pocos \textit{clusters}. El índice de Kakwani confirma la progresividad del programa, aunque decreciente entre regímenes; el índice de necesidad revela que los estados líderes en valor per cápita están subatendidos frente a la intensidad local de la pobreza. El artículo publica el \textit{pipeline}, los datos y las pruebas automatizadas, permitiendo su reproducción integral.

\vspace{6pt}\noindent\textbf{Palabras clave:} Bolsa Família; transferencias monetarias condicionadas; evaluación de políticas públicas; equidad distributiva; reproducibilidad.
\end{singlespace}

\clearpage

\section{Introdução}\label{sec:intro}


O Programa Bolsa Família (PBF), criado pela Lei n.\,10.836, de 9 de
janeiro de 2004, consolidou-se como o principal instrumento brasileiro
de transferência condicionada de renda. Em duas décadas, atravessou
três regimes institucionais: PBF clássico (2003\textendash 2021),
Auxílio Brasil (Nov.\,2021\textendash Fev.\,2023, instituído pela
Medida Provisória n.\,1.061/2021 sob o governo Bolsonaro) e Novo
Bolsa Família (Mar.\,2023 em diante, restabelecido pela Lei
n.\,14.601/2023 sob o governo Lula). Cada regime trouxe alterações
no desenho do benefício, na cobertura, no valor pago e nos critérios
de elegibilidade.

A literatura acadêmica produziu evidências consistentes sobre os
efeitos do PBF na redução da pobreza extrema, na escolarização
infantil, na frequência de pré-natal e na ampliação do consumo nas
regiões mais pobres do país (Soares, 2010; Paes-Sousa et al., 2011; Campello \& Neri, 2013). A literatura internacional reconheceu o
programa como referência mundial em transferência
condicionada (Fiszbein \& Schady, 2009), inspirando desenhos congêneres
em mais de quarenta países.

Apesar dessa centralidade, persistem três lacunas analíticas
distintas que este trabalho endereça em estrutura única.
\textbf{Primeiro}, o histórico consolidado dos microdados de
execução do programa permanece de acesso custoso: o Portal da
Transparência da CGU disponibiliza arquivos mensais (um por
competência) em CSV compactado, totalizando mais de 2,2 bilhões de
registros e cerca de 280\,GB descomprimidos no período
2013\textendash 2025. A consolidação, deflação e normalização
desses dados em painel UF\,$\times$\,Ano exige capacidade técnica
não trivial, o que limita o alcance da informação a um público
especializado. \textbf{Segundo}, a literatura empírica brasileira
sobre o PBF é majoritariamente descritiva ou ancorada em
\textit{surveys} (PNAD-C, POF) \textemdash{} não há estudo
quasi-experimental sobre as transições institucionais MP 1.061/2021
e Lei 14.601/2023, ainda que ambas constituam choques
identificacionais com cutoffs precisos. \textbf{Terceiro}, o debate
público sobre sustentabilidade fiscal do programa em \rs 140\,bilhões
anuais carece de aparato comparativo internacional sistemático.

Este artigo organiza-se em sete seções. A Seção
\ref{sec:teoria} revisa a evolução institucional. A Seção
\ref{sec:metodo} apresenta a metodologia. A Seção
\ref{sec:resultados} expõe os resultados descritivos. A Seção
\ref{sec:causal} apresenta a identificação causal das transições
institucionais. A Seção \ref{sec:equidade} apresenta os índices
formais de equidade. A Seção \ref{sec:benchmark} apresenta o
benchmark internacional. As Seções \ref{sec:discussao} e
\ref{sec:conclusao} discutem implicações e limites.


\section{Histórico Institucional}\label{sec:teoria}


\subsection{O Bolsa Família clássico (2003\textendash 2021)}

Criado pela Lei n.\,10.836/2004, o PBF unificou em um único arranjo
quatro programas anteriores (Bolsa Escola, Bolsa Alimentação,
Auxílio Gás e Cartão Alimentação) sob coordenação do então
Ministério do Desenvolvimento Social. O critério de elegibilidade
central foi a renda familiar per capita: famílias em pobreza
extrema (renda < \rs 89/mês até 2018) recebiam o benefício básico
acrescido de variáveis por composição familiar; famílias em pobreza
(renda entre \rs 89 e \rs 178) recebiam apenas as parcelas
variáveis. Ao longo de quase duas décadas, o programa manteve
cobertura entre 14 e 17 milhões de famílias, com valor médio anual
em \rs 35\textendash 41\,bilhões em valores reais (preços
dez/2021).

\subsection{O Auxílio Brasil (Nov.\,2021\textendash Fev.\,2023)}

A Medida Provisória n.\,1.061, de 9 de agosto de 2021,
posteriormente convertida na Lei n.\,14.284/2021, extinguiu o PBF e
criou o Auxílio Brasil em seu lugar. As principais alterações:
(i) elevação do piso para \rs 400/família, posteriormente
\rs 600/família em meio à campanha eleitoral de 2022;
(ii) ampliação da cobertura para aproximadamente 24 milhões de
famílias mediante incorporação de beneficiários do Cadastro Único
previamente fora da fila; (iii) modificação dos critérios de
condicionalidade. O período é caracterizado por forte politização e
vinculação ao calendário eleitoral.

\subsection{O Novo Bolsa Família (Mar.\,2023\textendash atual)}

A Lei n.\,14.601, de 19 de junho de 2023, revogou o Auxílio Brasil e
restabeleceu o Bolsa Família com nova configuração: (i) piso de
\rs 600/família mais adicionais por composição familiar (crianças,
gestantes, nutrizes); (ii) restabelecimento das condicionalidades
de educação e saúde; (iii) ampliação da cobertura, com inclusão de
famílias unipessoais antes restritas; (iv) integração com
Cadastro Único modernizado.

\subsection{Os marcos institucionais visualizados}

A Figura \ref{fig:timeline} sintetiza os marcos discutidos com as
faixas temporais correspondentes a cada regime.

\begin{figure}[htbp]
\centering
\includegraphics[width=0.96\linewidth]{fig01-timeline-pbf}
\caption{Linha do tempo institucional do Bolsa Família, Auxílio Brasil e Novo Bolsa Família (2003\textendash 2025)}
\label{fig:timeline}
\par\noindent{\small Fonte: Lei 10.836/2004, MP 1.061/2021, Lei 14.284/2021, Lei 14.601/2023.}
\end{figure}


\section{Metodologia}\label{sec:metodo}


\subsection{Fontes de dados}

Três fontes oficiais públicas: (i) microdados mensais do Portal da
Transparência da CGU (CGU, 2025), abrangendo PBF
(2013\textendash 2021), Auxílio Brasil (2021\textendash 2023) e
Novo Bolsa Família (2023\textendash atual); (ii) IBGE/SIDRA Tabela
6579 (população residente estimada por UF) (IBGE, 2024);
(iii) Banco Central do Brasil, série SGS 433 (IPCA mensal para
deflação base dezembro/2021) (Banco Central do Brasil, 2024). Para o índice
de Kakwani, IDH-M PNUD/Atlas Brasil 2010 (último Censo
completo) (PNUD et al., 2013). Para a construção do
tratamento da identificação causal, taxa de pobreza monetária
PNAD-C 2019 (linha US\$~5,50/dia PPP, IBGE Tabela
6688) (IBGE, 2019). Para o benchmark internacional,
\textit{World Bank ASPIRE} (World Bank, 2024a), CEPAL
\textit{Panorama Social} (CEPAL, 2024) e fontes nacionais por
programa (ANSES, SEDESOL, DPS, MEFP\,/\,UDAPE).

\subsection{Pipeline de dados}

O pipeline adota arquitetura medallion (Armbrust et al., 2021; Kukreja, 2021) implementada em Apache
Spark (Zaharia et al., 2016) sobre Databricks Free Edition com
persistência em Delta Lake.
\textbf{Bronze} ingere os ZIPs mensais via Auto Loader (Spark
Structured Streaming) preservando o registro original em
\textit{string-only} \textemdash{} padrão da plataforma para evitar
divergências entre o que está no arquivo (latin-1, separador
decimal vírgula) e o esquema Delta inferido.
\textbf{Silver} aplica deduplicação de beneficiários por NIS,
correção do swap \texttt{mes\_competencia}/\texttt{mes\_referencia}
presente nos arquivos de novembro/2021 (transição PBF\,$\to$\,AB),
parsing dos valores monetários no formato brasileiro e agregação
mensal por (UF, Ano, Mês).
\textbf{Gold} consolida em painel UF\,$\times$\,Ano com deflação
IPCA, normalização per capita (IBGE) e cálculo dos indicadores
derivados, filtrando anos não-completos pelo invariante
\texttt{n\_months == 12} (com \textit{logging} dos descartados).

\begin{tcolorbox}[
  colback=blue!2,colframe=blue!60!black,
  boxrule=0.6pt,arc=2pt,
  left=8pt,right=8pt,top=4pt,bottom=4pt,
  title={\small\textbf{Nota metodológica — tratamento da transição PBF\,$\to$\,AB em novembro de 2021}},
  fonttitle=\small\bfseries,
  ]\small
O mês de novembro de 2021 contém registros simultâneos de PBF e
Auxílio Brasil, reflexo da transição administrativa. O silver
sintetiza um identificador composto \texttt{PBF\_AUX\_SUM} que
\textit{soma} as parcelas das duas modalidades para beneficiários
presentes em ambos, evitando dupla contagem; quando a coluna
\texttt{mes\_competencia} aparece com valor distinto de
\texttt{mes\_referencia} (artefato dos arquivos da CGU nesse mês),
o silver normaliza pelo \texttt{mes\_competencia}. Detalhes de
implementação: \texttt{pipelines/notebooks/silver/pbf\_total\_uf\_mes.py}.
\end{tcolorbox}

\subsection{Indicadores construídos}

Para cada par (UF, Ano): (i) número de beneficiários distintos
(\textit{deduplicação por NIS});
(ii) valor pago nominal e em \rs 2021;
(iii) per beneficiário (valor pago / beneficiários);
(iv) per capita (valor pago / população);
(v) penetração (beneficiários / população);
(vi) coeficiente de variação per capita entre UFs por exercício.

\subsection{Identificação causal}

Sejam $y_{u,t}$ o valor anual pago em \rs 2021 per capita na UF $u$
no ano $t$. \textbf{Tratamento.} $T_u = 1$ se UF $u$ está no quartil
superior do \textit{déficit de cobertura} pré-choque, definido como
$d_u = \mathit{pobreza}_u^{\,\mathrm{PNAD-C\,2019}} - \overline{\mathit{penetra\c{c}\~ao}}_u^{\,2018\text{-}2020}$.
A intuição é que UFs com pobreza alta e penetração baixa têm maior
folga para absorver a expansão de cobertura do Auxílio Brasil
\textemdash{} são as mais \textit{tratadas} pela mudança institucional
de novembro/2021. \textbf{Modelos.}
(i)~DiD~2$\times$2 sobre a diferença
$\Delta y_u = \overline{y_{u,t}}_{t \in \{2022,...,2025\}} - \overline{y_{u,t}}_{t \in \{2018,2019,2020\}}$:
$\Delta y_u = \alpha + \beta\,T_u + \varepsilon_u$, $\hat{\beta}$
HC3.
(ii)~TWFE com efeitos fixos de UF e ano e erros clusterizados por
UF: $y_{u,t} = \alpha + \beta\,(\mathit{Post}_t \cdot T_u)
+ \gamma_u + \delta_t + \varepsilon_{u,t}$, restringido a
$2018 \leq t \leq 2025$.
(iii)~\textit{Wild-cluster bootstrap} com $999$ réplicas Rademacher
por UF, dada a inconsistência conhecida de erros clusterizados
assintóticos quando o número de \textit{clusters} é
pequeno (Cameron et al., 2008): o $p$-valor robusto é o
quantil empírico de $|t^{(b)}| \geq |t^{\mathrm{obs}}|$ sobre as
réplicas.
\textbf{Salvaguardas adicionais.}
(iv)~Replicação cross-shock sobre Lei 14.601/2023, com
$T_u$ idêntico mas pré $\in \{2021,2022\}$ e pós $\in \{2023,2024,2025\}$.
(v)~Teste informal de \textit{parallel trends} pré-tratamento via
regressão $y_{u,t} = \alpha + \rho\,t + \pi\,T_u + \tau\,(T_u \cdot t)
+ \varepsilon$ sobre $t \in \{2018,2019,2020\}$, com clustering por
UF; a hipótese de tendências paralelas pré-tratamento corresponde
a $\tau = 0$ (Roth, 2022).
(vi)~Placebo: substituir o \textit{cutoff} para 2020 (sem mudança
institucional) e re-estimar.
(vii)~\textit{Leave-one-out} retirando cada UF.

\subsection{Índices formais de equidade}

\textbf{Índice de Kakwani} (Kakwani, 1977): $K = C_y - G_z$, onde
$C_y$ é o índice de concentração da variável $y$ ordenada pelo
\textit{ranking} $z$ e $G_z$ é o Gini do próprio \textit{ranking}.
Ordenamos as 27 UFs pelo IDH-M crescente. $K < 0$ indica
\textit{progressividade} (mais alocação a UFs de menor IDH-M);
$K > 0$ indica \textit{regressividade}; $K \approx 0$ indica
neutralidade. Calculamos $K_{\mathrm{PBF}}$ em 2018, 2020 (pré-AB),
2022 (AB), 2024 (NBF) e $K_{\mathrm{emendas}}$ nos mesmos
exercícios, com IC~95\,\% \textit{bootstrap} (1000 réplicas).
\textbf{Índice de necessidade.}
$N_u = \mathit{pobreza}_u \cdot \mathit{pop}_u\,/\,\sum_v
\mathit{pobreza}_v \cdot \mathit{pop}_v$.
A razão $R_u = \mathit{share}_{\mathrm{PBF},u} / N_u$ mede
adequação da alocação à necessidade local: $R > 1$ indica
sobre-cobertura em relação à intensidade da pobreza local;
$R < 1$ indica sub-cobertura.

\subsection{Benchmark internacional}

Comparação direta entre o NBF brasileiro e quatro CCTs
latino-americanos: AUH (Argentina, 2009\textendash), Prospera
(México, 1997\textendash 2019, extinto pelo governo
AMLO), Más Familias en Acción (Colômbia, 2001\textendash) e
Renta Dignidad (Bolívia, 2008\textendash, universalmente vinculado
à idade). Métricas: cobertura (\% da população atendida), custo
(\% do PIB anual), valor per beneficiário em US\$~PPP~2021
(World Bank ICP, 2024) e custo administrativo declarado
(\% do gasto total).

\subsection{Limitações}

(i)~A unidade de análise é a \textit{família} beneficiária, não o
indivíduo \textemdash{} o impacto sobre o orçamento familiar pode
diferir do per capita populacional;
(ii)~o exercício 2026 está em curso e foi excluído da análise (o
gold filtra por \texttt{n\_months == 12});
(iii)~o Auxílio Emergencial (2020\textendash 2021), benefício
distinto, não é contemplado;
(iv)~a identificação causal sobre MP 1.061/2021 está sujeita a
três limites declarados: número pequeno de \textit{clusters}
(N=27), coexistência com choques contemporâneos (recuperação
pós-COVID, inflação) e rejeição do teste de \textit{parallel trends}
pré-tratamento, conforme detalhado na Seção
\ref{sec:causal:robustez};
(v)~o IDH-M utilizado refere-se ao Censo de 2010 (último
disponível para o índice na série completa por UF);
(vi)~a taxa de pobreza PNAD-C 2019 é \textit{snapshot}
pré-pandêmico e pode subestimar o tratamento em UFs onde a
pobreza ampliou em 2020\textendash 2021;
(vii)~o benchmark internacional usa anos próximos mas não
idênticos por país, limitando comparabilidade exata.


\section{Resultados Descritivos}\label{sec:resultados}


\subsection{Evolução temporal: três regimes empíricos}

A Tabela \ref{tab:evol-pbf} apresenta a série temporal nacional
consolidada. Observam-se claramente os dois saltos: entre 2021 e
2022 (Auxílio Brasil) com fator real $2,16\times$, e entre 2022 e
2023 (Novo Bolsa Família) com fator real $2,07\times$. As
variações de magnitude $> 30\,\%$ no per beneficiário em 2022
($+36\,\%$) e 2023 ($+102\,\%$) indicam transformação institucional,
não apenas atualização monetária.

\begin{table}[H]
\centering
\caption{Bolsa Família/AB/NBF: beneficiários, valor pago e per capita Brasil, 2013\textendash 2025}
\begin{tabular}{rrrrrr}
\toprule
\textbf{Ano} & \textbf{Beneficiários} & \textbf{Pago real} & \textbf{Pago nominal} & \textbf{Per beneficiário} & \textbf{Per capita} \\
   & (milhões)             & (\rs bi, 2021) & (\rs bi, correntes) & (\rs/ano, 2021) & (\rs/hab, 2021) \\
\midrule
2013 & 16,43 &  39,92 &  24,89 & 2.429 & 198,60 \\
2014 & 15,78 &  40,99 &  27,19 & 2.598 & 202,14 \\
2015 & 16,16 &  37,66 &  27,65 & 2.330 & 184,21 \\
2016 & 16,02 &  36,53 &  28,51 & 2.280 & 177,26 \\
2017 & 16,13 &  36,16 &  29,05 & 2.241 & 174,11 \\
2018 & 16,99 &  36,75 &  30,63 & 2.163 & 176,25 \\
2019 & 14,99 &  35,84 &  31,16 & 2.391 & 170,56 \\
2020 & 14,75 &  35,22 &  32,00 & 2.388 & 166,33 \\
2021 & 15,02 &  30,39 &  30,39 & 2.023 & 142,44 \\
2022 & 23,85 &  65,76 &  69,56 & 2.757 & 308,61 \\
2023 & 24,51 & 136,27 & 150,82 & 5.560 & 640,29 \\
2024 & 23,18 & 140,88 & 163,45 & 6.078 & 662,70 \\
2025 & 22,29 & 129,87 & 157,10 & 5.826 & 608,51 \\
\bottomrule
\end{tabular}
\par\noindent{\small Fonte: elaboração própria, microdados CGU, IPCA-BCB para deflação base dez/2021. Dados gerados pelo gold dataset \texttt{gold\_pbf\_estados\_df.json} (\COMPILESHA).}
\label{tab:evol-pbf}
\end{table}

A Figura \ref{fig:evolution} apresenta visualmente os mesmos dados
em eixo duplo, destacando os dois saltos.

\begin{figure}[htbp]
\centering
\includegraphics[width=0.92\linewidth]{fig02-evolution-dual}
\caption{Evolução do valor pago anual (R\$ bi, 2021) e do número de beneficiários (milhões), 2013\textendash 2025}
\label{fig:evolution}
\par\noindent{\small Fonte: elaboração própria, microdados CGU.}
\end{figure}

\subsection{Per beneficiário: a transformação qualitativa do programa}

A Figura \ref{fig:per-benef} apresenta a evolução do valor médio
pago por beneficiário em valores reais. O patamar histórico de
\rs 2.000\textendash 2.600/ano salta para \rs 5.500\textendash
6.000/ano a partir de 2023, em virtude da elevação do piso e da
restruturação dos adicionais por composição familiar. A Figura
\ref{fig:yoy} apresenta a variação ano-a-ano (YoY).

\begin{figure}[htbp]
\centering
\includegraphics[width=0.92\linewidth]{fig03-per-benef}
\caption{Valor anual pago por beneficiário (R\$ 2021), 2013\textendash 2025, com faixas temporais dos três regimes}
\label{fig:per-benef}
\par\noindent{\small Fonte: elaboração própria. Faixas verticais identificam PBF clássico (verde), Auxílio Brasil (vermelho) e Novo Bolsa Família (azul).}
\end{figure}

\begin{figure}[htbp]
\centering
\includegraphics[width=0.92\linewidth]{fig12-yoy-growth}
\caption{Variação ano-a-ano (\%) do valor pago por beneficiário, 2014\textendash 2025}
\label{fig:yoy}
\par\noindent{\small Fonte: elaboração própria.}
\end{figure}

\subsection{Distribuição regional}

A Figura \ref{fig:regional} apresenta a distribuição entre as cinco
regiões brasileiras em três métricas \textemdash{} per capita
ponderado pela população, valor agregado em \rs bi e penetração
(\% da população beneficiária). Em 2025, o
\textbf{Nordeste} domina nas três métricas (\rs 1.051/hab,
\rs 60,2\,bi, penetração 17,6\,\%), seguido pelo \textbf{Norte}.
Sudeste e Centro-Oeste apresentam patamar intermediário; o
\textbf{Sul} apresenta os menores indicadores em todas as métricas
(\rs 283/hab, penetração 5,3\,\%) \textemdash{} razão
Nordeste/Sul de 3,72:1 em per capita ponderado.

\begin{figure}[htbp]
\centering
\includegraphics[width=\linewidth]{fig04-regional}
\caption{Bolsa Família por região: per capita ponderado, valor absoluto e penetração (2025)}
\label{fig:regional}
\par\noindent{\small Fonte: elaboração própria, microdados CGU + IBGE.}
\end{figure}

\subsection{Top-10 unidades federativas}

A Tabela \ref{tab:topabs-pbf} e a Figura \ref{fig:topabs}
apresentam as 10 UFs de maior volume acumulado em valores reais no
período. O ranking absoluto é dominado por estados de grande
população (Bahia, São Paulo, Pernambuco, Minas Gerais, Ceará). Como
documentado na próxima subseção, o ranking per capita inverte
substancialmente a ordem.

\begin{table}[H]
\centering
\caption{Top-10 UFs por valor pago acumulado em PBF/AB/NBF em valores reais (\rs bi, 2021), 2013\textendash 2025}
\begin{tabular}{rlr}
\toprule
\textbf{Posição} & \textbf{UF} & \textbf{Pago acumulado (\rs bi, 2021)} \\
\midrule
1 & Bahia            & 99,35 \\
2 & São Paulo        & 87,63 \\
3 & Pernambuco       & 63,10 \\
4 & Minas Gerais     & 60,17 \\
5 & Ceará            & 58,55 \\
6 & Rio de Janeiro   & 56,71 \\
7 & Maranhão         & 54,93 \\
8 & Pará             & 54,92 \\
9 & Paraíba          & 28,45 \\
10 & Amazonas        & 26,54 \\
\bottomrule
\end{tabular}
\par\noindent{\small Fonte: elaboração própria a partir do gold dataset \texttt{gold\_pbf\_estados\_df.json}, microdados CGU. Valores agregados como soma de \texttt{valor\_2021} (deflacionado dez/2021) sobre o período. Notar a inversão: SP, segundo no ranking absoluto, cai à 25ª posição no ranking per capita (Tabela \ref{tab:percapita-pbf}).}
\label{tab:topabs-pbf}
\end{table}

\begin{figure}[htbp]
\centering
\includegraphics[width=0.85\linewidth]{fig05-top10-abs}
\caption{Top-10 UFs em valor pago acumulado em PBF/AB/NBF (R\$ bi, 2021), 2013\textendash 2025}
\label{fig:topabs}
\par\noindent{\small Fonte: elaboração própria.}
\end{figure}

\subsection{Análise per capita: focalização territorial}

A Tabela \ref{tab:percapita-pbf} e a Figura \ref{fig:rank}
apresentam o ranking per capita 2025. Maranhão lidera com
\rs 1.168,51/hab (\textbf{18,8\,\%} da população beneficiária);
Santa Catarina é a última com \rs 167,35/hab (\textbf{3,3\,\%} de
penetração). A razão entre os extremos é de
\textbf{6,98:1}\,\textemdash{} padrão moderado quando comparado às
emendas parlamentares (vide Seção \ref{sec:cv}).

\begin{table}[H]
\centering
\caption{Top-10 e bottom-5 UFs por valor pago per capita do PBF/AB/NBF (\rs/hab, 2021), 2025}
\begin{tabular}{rlrr}
\toprule
\textbf{Pos.} & \textbf{UF} & \textbf{Per capita (\rs/hab)} & \textbf{Penetração (\%)} \\
\midrule
1 & Maranhão              & \textbf{1.168,51} & 18,8 \\
2 & Piauí                 & 1.119,81 & 18,6 \\
3 & Alagoas               & 1.077,29 & 17,8 \\
4 & Bahia                 & 1.052,09 & 17,8 \\
5 & Pernambuco            & 1.050,98 & 17,7 \\
6 & Sergipe               & 1.047,55 & 16,5 \\
7 & Amapá                 & 1.041,57 & 18,0 \\
8 & Amazonas              & 1.039,64 & 16,2 \\
9 & Acre                  & 1.028,75 & 18,1 \\
10 & Pará                 & 1.027,26 & 17,2 \\
\multicolumn{4}{c}{\dots\ (UFs intermediárias omitidas) \dots} \\
23 & Distrito Federal     & 356,85 &  6,3 \\
24 & Rio Grande do Sul    & 337,62 &  6,2 \\
25 & São Paulo            & 325,33 &  6,0 \\
26 & Paraná               & 309,86 &  5,7 \\
27 & Santa Catarina       & \textbf{167,35} & 3,3 \\
\bottomrule
\end{tabular}
\par\noindent{\small Fonte: elaboração própria, microdados CGU + IBGE/SIDRA Tabela 6579 (população residente estimada). Razão Maranhão/Santa Catarina = 6,98:1.}
\label{tab:percapita-pbf}
\end{table}

\begin{figure}[htbp]
\centering
\includegraphics[width=0.85\linewidth]{fig06-choropleth-percapita}
\caption{Mapa coroplético — Bolsa Família per capita por UF (R\$/hab, 2021), 2025}
\label{fig:choro-pc}
\par\noindent{\small Fonte: elaboração própria, GeoJSON IBGE + microdados CGU.}
\end{figure}

\begin{figure}[htbp]
\centering
\includegraphics[width=0.85\linewidth]{fig07-choropleth-absoluto}
\caption{Mapa coroplético — Bolsa Família valor pago por UF (\rs bi, 2021), 2025}
\label{fig:choro-abs}
\par\noindent{\small Fonte: elaboração própria. SP, MG, BA dominam em valores absolutos; o Nordeste domina em per capita.}
\end{figure}

\begin{figure}[htbp]
\centering
\includegraphics[width=0.7\linewidth]{fig08-heatmap}
\caption{Heatmap UF $\times$ Ano — Bolsa Família per capita (R\$/hab, 2021), 2018\textendash 2025}
\label{fig:heatmap}
\par\noindent{\small Fonte: elaboração própria. Linhas verticais tracejadas marcam Auxílio Brasil (Nov/2021) e NBF (Mar/2023). UFs ordenadas em ordem decrescente do per capita 2025.}
\end{figure}

\begin{figure}[htbp]
\centering
\includegraphics[width=0.7\linewidth]{fig09-per-capita-ranking}
\caption{Ranking horizontal: Bolsa Família per capita por UF, 2025}
\label{fig:rank}
\par\noindent{\small Fonte: elaboração própria.}
\end{figure}

\subsection{Penetração e per capita: escopo da associação}

A Figura \ref{fig:pen} apresenta a relação entre penetração
($\%$ da população beneficiária) e per capita pago em 2025. A
correlação é positiva ($\rho \approx 0,98$, IC 95\,\% bootstrap
$[0,97; 0,99]$, $N=27$ UFs, 1000 réplicas), em parte por
construção parcial: as duas variáveis compartilham o numerador
(\texttt{n\_benef} \,/\, \texttt{populacao} versus \texttt{valor}
\,/\, \texttt{populacao}). A magnitude próxima de 1 reflete
homogeneidade do valor médio pago por família entre UFs após o NBF.

\begin{figure}[htbp]
\centering
\includegraphics[width=0.92\linewidth]{fig10-penetracao}
\caption{Penetração ($\%$ da população beneficiária) vs per capita 2025 por UF, com $\rho$ e IC 95\,\% bootstrap}
\label{fig:pen}
\par\noindent{\small Fonte: elaboração própria.}
\end{figure}

\subsection{Coeficiente de variação: comparação descritiva com Emendas
Parlamentares}\label{sec:cv}

A Figura \ref{fig:cv} sobrepõe o CV per capita do PBF e o do
programa de emendas parlamentares federais. O CV do PBF oscila
entre $0,428$ e $0,571$ ao longo do período; o de emendas, entre
$0,570$ e $1,776$ (com pico em 2017, primeiro ano após a EC 86 que
tornou o RP6 obrigatório \textemdash{} a anotação na figura
detalha o transitório).

A diferença entre os dois CVs é descritiva: a Seção
\ref{sec:equidade} substitui esse contraste descritivo por
comparação cardinal via índice de Kakwani.

\begin{figure}[htbp]
\centering
\includegraphics[width=0.92\linewidth]{fig11-cv-comparison}
\caption{Coeficiente de variação per capita por UF — Bolsa Família vs Emendas Parlamentares, 2013\textendash 2025, com anotação no outlier 2017 das emendas}
\label{fig:cv}
\par\noindent{\small Fonte: elaboração própria, microdados CGU + IBGE.}
\end{figure}


\section{Identificação causal das transições institucionais}\label{sec:causal}


A Seção \ref{sec:resultados} estabeleceu, descritivamente, que o
programa atravessou dois pontos de mudança quase simultâneos: MP
1.061/2021 (Auxílio Brasil) e Lei 14.601/2023 (Novo Bolsa Família).
Esta seção avalia, em desenho quasi-experimental, se o salto
observado em valor per capita corresponde a efeito causal
identificável ou se reflete tendências pré-existentes.

\subsection{Construção do tratamento e amostra}

O tratamento $T_u$ classifica UFs pelo \textit{déficit de cobertura}
pré-choque, conforme Seção \ref{sec:metodo}. O \textit{cutoff} no
quartil superior produz \textbf{7 UFs tratadas} (AL, AM, CE, MA,
PA, PB, PI) e 20 UFs de controle. A Figura \ref{fig:barbell}
visualiza a transição pré $\to$ pós em per capita.

\begin{figure}[htbp]
\centering
\includegraphics[width=0.86\linewidth]{fig13-causal-barbell}
\caption{\textit{Barbell} — Bolsa Família per capita por UF, pré (média 2018\textendash 2020) vs pós (média 2022\textendash 2025). Vermelho: UF tratada (Q4 do déficit pré-choque)}
\label{fig:barbell}
\par\noindent{\small Fonte: elaboração própria, microdados CGU + IBGE/PNAD-C.}
\end{figure}

\subsection{Resultados}

A Tabela \ref{tab:did-AB} apresenta os coeficientes do DiD~2$\times$2
e do TWFE para o choque MP 1.061/2021, com erros HC3 e clusterizados
por UF. A Tabela \ref{tab:did-NBF} apresenta a replicação cross-shock
sobre Lei 14.601/2023.

\begin{table}[H]
\centering
\caption{DiD sobre MP 1.061/2021 (pré 2018\textendash 2020 vs pós 2022\textendash 2025)}
\begin{tabular}{lrrrr}
\toprule
\textbf{Modelo} & $\hat{\beta}$ (R\$/hab) & SE & IC 95\,\% & $p$-valor \\
\midrule
DiD 2$\times$2 (HC3)              & $+223{,}75$ & 38{,}81 & $[+147{,}69;\,+299{,}82]$ & $< 0{,}001$ \\
TWFE clusterizado por UF (N=27)   & $+233{,}27$ & 42{,}00 & $[+150{,}96;\,+315{,}59]$ & $< 0{,}001$ \\
\textit{Wild-cluster bootstrap} (999 sims) & --- & --- & --- & $0{,}493$ \\
\bottomrule
\end{tabular}
\par\noindent{\small Fonte: \texttt{causal\_analysis\_pbf.py}. $\bar{\Delta}$ Treated $=$ \rs $+609{,}3$/hab; $\bar{\Delta}$ Control $=$ \rs $+385{,}6$/hab.}
\label{tab:did-AB}
\end{table}

\begin{table}[H]
\centering
\caption{DiD sobre Lei 14.601/2023 — replicação cross-shock (pré 2021\textendash 2022 vs pós 2023\textendash 2025)}
\begin{tabular}{lrrrr}
\toprule
\textbf{Modelo} & $\hat{\beta}$ (R\$/hab) & SE & IC 95\,\% & $p$-valor \\
\midrule
DiD 2$\times$2 (HC3) & $+271{,}05$ & 44{,}99 & $[+182{,}86;\,+359{,}23]$ & $< 0{,}001$ \\
\bottomrule
\end{tabular}
\par\noindent{\small Fonte: \texttt{causal\_analysis\_pbf.py}.}
\label{tab:did-NBF}
\end{table}

A Figura \ref{fig:event} mostra as trajetórias médias anuais de
Treated vs Control com marcadores verticais nos dois \textit{cutoffs}
institucionais.

\begin{figure}[htbp]
\centering
\includegraphics[width=0.92\linewidth]{fig14-event-study}
\caption{Trajetórias médias de per capita real (R\$/hab, 2021): Treated (Q4 déficit pré-choque) vs Control}
\label{fig:event}
\par\noindent{\small Fonte: elaboração própria. Linhas verticais marcam MP 1.061/2021 (Nov/2021) e Lei 14.601/2023 (Mar/2023).}
\end{figure}

\subsection{Robustez e limitações da identificação}\label{sec:causal:robustez}

\textbf{Wild-cluster bootstrap.} Sob HC3 e clustering assintótico,
o efeito é ordens de magnitude maior que o erro padrão; sob WCB com
999 réplicas Rademacher, o $p$-valor sobe para $0,493$. Esse
contraste é uma manifestação clássica da inconsistência de erros
clusterizados quando o número de \textit{clusters} é
pequeno (Cameron et al., 2008). \textbf{O efeito não pode ser
distinguido de zero sob a métrica robusta apropriada para N=27}.

\textbf{Parallel trends.} A regressão pré-tratamento
$y_{u,t} = \alpha + \rho t + \pi T_u + \tau (T_u \cdot t)
+ \varepsilon$ sobre $t \in \{2018,2019,2020\}$ produz
$\hat\tau = -9,03$ ($p < 0,001$, clustering por UF). A hipótese
de tendências paralelas é \textbf{rejeitada}: já havia divergência
sistemática entre Treated e Control antes de novembro de 2021. Sob
o \textit{ônus} de Roth (2022), parte do efeito estimado pode
refletir essa divergência pré-existente.

\textbf{Placebo.} O DiD com \textit{cutoff} em 2020 (sem mudança
institucional) retorna $\hat\beta = -11,72$ ($p = 0,004$),
sugerindo presença de tendência diferencial não-relacionada a MP
1.061/2021.

\textbf{Leave-one-out.} A retirada sequencial de cada UF resulta
em $\hat\beta \in [+209,70;\,+237,56]$ (média $+223,75$, dp $7,47$),
indicando estabilidade ao subconjunto amostral.

\textbf{Interpretação.} Os pontos (i) WCB $p = 0,49$ e (ii)
parallel-trends rejeitado são \textit{red flags} identificacionais.
A magnitude do salto descritivo é incontroversa; a atribuição
\textit{causal} da diferença a MP 1.061/2021 \textbf{não é sustentada
pela métrica robusta apropriada para N=27 \textit{clusters}}. O
desenho aqui apresentado é o melhor possível com os dados
disponíveis em frequência anual e cobertura por UF \textemdash{} a
identificação rigorosa do efeito requer microdados mensais por
município com Cadastro Único, que não estão publicamente disponíveis
em série completa.


\section{Análise formal de equidade}\label{sec:equidade}


\subsection{Índice de Kakwani sobre PBF e emendas, 2018\textendash 2024}

A Tabela \ref{tab:kakwani} apresenta $K = C_y - G_z$ para PBF
per capita ordenado por IDH-M e para emendas per capita no mesmo
\textit{ranking}, em quatro exercícios. Os intervalos de confiança
\textit{bootstrap} (1000 réplicas) acompanham cada estimativa.

\begin{table}[H]
\centering
\caption{Kakwani $K = C_y - G_{\mathrm{IDHM}}$ — PBF e emendas, com IC 95\,\% bootstrap}
\begin{tabular}{rlrrrr}
\toprule
\textbf{Ano} & \textbf{Programa/Regime} & $K$ & IC inf. & IC sup. & Sinal \\
\midrule
2018 & PBF clássico        & $-0{,}3340$ & $-0{,}3942$ & $-0{,}2474$ & progressivo \\
2020 & PBF pré-AB          & $-0{,}3238$ & $-0{,}3810$ & $-0{,}2395$ & progressivo \\
2022 & Auxílio Brasil      & $-0{,}2749$ & $-0{,}3229$ & $-0{,}2045$ & progressivo \\
2024 & Novo Bolsa Família  & $-0{,}2582$ & $-0{,}3039$ & $-0{,}1898$ & progressivo \\
\midrule
2018 & Emendas pré-AB      & $-0{,}1734$ & $-0{,}3096$ & $-0{,}0263$ & progressivo (fraco) \\
2022 & Emendas 2022        & $-0{,}1312$ & $-0{,}2207$ & $-0{,}0246$ & progressivo (fraco) \\
2024 & Emendas 2024        & $-0{,}1548$ & $-0{,}2701$ & $-0{,}0251$ & progressivo (fraco) \\
\bottomrule
\end{tabular}
\par\noindent{\small Fonte: \texttt{equity\_indices\_pbf.py}. UFs ordenadas por IDH-M (PNUD/Atlas Brasil 2010) crescente. Bootstrap 1000 réplicas.}
\label{tab:kakwani}
\end{table}

Dois resultados distintos da literatura descritiva pré-existente:
(i) PBF/NBF é consistentemente mais progressivo que emendas
\textemdash{} a magnitude é cardinal, não apenas comparativa de CVs;
(ii) \textbf{a progressividade do programa caiu monotonicamente
entre 2018 e 2024} ($-0,334 \to -0,258$), sugerindo que a expansão
de cobertura via Auxílio Brasil e NBF incorporou famílias em UFs
de IDH-M mais alto na proporção, atenuando o gradiente
progressivo. O programa ficou mais grande mas, por essa métrica,
ligeiramente menos focado nos estados com pior IDH-M.

A Figura \ref{fig:kakwani} apresenta as curvas de concentração de
PBF/NBF (2024) e emendas (2024) sobre as UFs ordenadas por IDH-M.

\begin{figure}[htbp]
\centering
\includegraphics[width=0.92\linewidth]{fig15-kakwani-curve}
\caption{Curvas de concentração: PBF/NBF (2024) vs Emendas Parlamentares (2024), UFs ordenadas por IDH-M crescente}
\label{fig:kakwani}
\par\noindent{\small Fonte: elaboração própria.}
\end{figure}

\subsection{Índice de necessidade — adequação à intensidade da pobreza local}

A Figura \ref{fig:need} apresenta a razão
$R_u = \mathit{share}_{\mathrm{PBF},u} / \mathit{share}_{\mathit{necessidade},u}$
em 2024, ordenada do menor ao maior. UFs com $R > 1$ recebem
proporcionalmente mais que sua \textit{fatia de necessidade}
(definida como pobreza\,$\times$\,população); UFs com $R < 1$ estão
sub-cobertas em relação à intensidade da pobreza local.

\begin{figure}[htbp]
\centering
\includegraphics[width=0.85\linewidth]{fig16-need-vs-coverage}
\caption{Razão $R = \mathit{share}_{\mathrm{PBF}} / \mathit{share}_{\mathit{necessidade}}$ por UF, 2024. Linha tracejada em $R = 1$}
\label{fig:need}
\par\noindent{\small Fonte: elaboração própria, microdados CGU + PNAD-C 2019.}
\end{figure}

Resultado contraintuitivo distinto da análise per capita simples:
\textbf{UFs aparentemente \textit{campeãs} em per capita (Maranhão,
Roraima, Rondônia) estão de fato sub-atendidas em relação à
intensidade da pobreza local} ($R = 0,89; 0,80; 0,84$). Isso
ocorre porque a \textit{intensidade} da pobreza nessas UFs é tão
elevada que mesmo o per capita líder não compensa o número de
famílias em situação de pobreza monetária. UFs ricas como Distrito
Federal, Rio de Janeiro e Rio Grande do Sul aparecem com $R > 1$
\textemdash{} a pobreza nelas é menor em proporção, e a presença de
estruturas urbanas e de Cadastro Único mais consolidadas favorecem
sobre-cobertura relativa.


\section{Benchmark internacional}\label{sec:benchmark}


A Tabela \ref{tab:benchmark} apresenta o NBF brasileiro lado a lado
com quatro CCTs latino-americanos. A Figura \ref{fig:benchmark}
visualiza simultaneamente cobertura, custo (\% PIB) e per
beneficiário em US\$~PPP~2021.

\begin{table}[H]
\centering
\caption{CCTs na América Latina: cobertura, custo e per beneficiário}
\begin{tabular}{lrrrrr}
\toprule
\textbf{País / Programa} & \textbf{Ano} & \textbf{Cobert.\ (\% pop)} & \textbf{Custo (\% PIB)} & \textbf{Per benef.\ (US\$ PPP/ano)} & \textbf{Famílias (mi)} \\
\midrule
Brasil / NBF              & 2024 & 33,1\,\% & 1,10\,\% & 2.179 & 23,18 \\
Argentina / AUH           & 2023 &  9,5\,\% & 0,55\,\% & 1.380 &  2,20 \\
México / Prospera         & 2018 & 27,0\,\% & 0,30\,\% &   720 &  6,50 \\
Colômbia / MFA            & 2023 & 14,5\,\% & 0,40\,\% & 1.100 &  2,60 \\
Bolívia / Renta Dignidad  & 2023 & 11,2\,\% & 1,20\,\% & 1.670 &  1,40 \\
\bottomrule
\end{tabular}
\par\noindent{\small Fontes: World Bank ASPIRE database (World Bank, 2024a), CEPAL Panorama Social 2024 (CEPAL, 2024), MDS/CGU (Brasil), ANSES (Argentina), SEDESOL (México), DPS (Colômbia), MEFP (Bolívia). PPP 2021 da ICP do World Bank (World Bank, 2024b).}
\label{tab:benchmark}
\end{table}

\begin{figure}[htbp]
\centering
\includegraphics[width=\linewidth]{fig17-cct-international}
\caption{Comparativo CCT — Brasil vs Argentina, México, Colômbia e Bolívia: cobertura, custo (\%PIB) e per beneficiário US\$ PPP 2021}
\label{fig:benchmark}
\par\noindent{\small Fonte: elaboração própria.}
\end{figure}

Quatro observações comparativas. \textbf{(i) Cobertura.} O NBF
brasileiro a $\sim$\,33\,\% da população está em patamar superior à
maioria dos CCTs latino-americanos, apenas comparável ao pico do
Prospera mexicano em sua fase pré-extinção (2018). \textbf{(ii)
Custo (\%PIB).} Brasil em 1,10\,\% é elevado para o padrão regional
(Argentina 0,55\,\%, México 0,30\,\%, Colômbia 0,40\,\%) \textemdash{}
superado apenas pela Bolívia (1,20\,\%, mas universal por idade,
não focalizado por renda). \textbf{(iii) Per beneficiário.} O NBF
entrega US\$~PPP~2.179/família/ano, \textbf{o maior valor da
região} \textemdash{} $\sim$\,58\,\% acima da Argentina, três vezes
o México e duas vezes a Colômbia. \textbf{(iv) Custo
administrativo.} Brasil reporta o menor custo administrativo
declarado ($\sim$\,1,5\,\%), reflexo da operação bancária
centralizada na CEF e do CadÚnico amortizado.


\section{Aprendizado de máquina: o Bolsa Família como \textit{proxy} de pobreza e desenvolvimento}\label{sec:ml}


Uma literatura consolidada emprega a cobertura do Bolsa Família como
\textit{proxy} territorial de pobreza e de nível de desenvolvimento
local, sobretudo onde indicadores diretos e atualizados são escassos na
granularidade desejada (Rasella et al., 2013; Hoffmann, 2006; Rocha, 2008). Estudos epidemiológicos e econômicos usam a penetração
do programa como variável de exposição associada a desfechos sociais
(p.\,ex.\ mortalidade infantil ao nível municipal) (Rasella et al., 2013).
Esta seção testa \textbf{explicitamente}, com aprendizado de máquina, em
que medida indicadores \textit{internos} do programa recuperam medidas
consagradas de pobreza (taxa de pobreza monetária PNAD-C 2019) e de
desenvolvimento (IDH-M 2010). A exposição é escalonada de propósito: do
modelo mais simples e interpretável (regressão) ao mais flexível
(florestas aleatórias).

\subsection{Dados, variáveis e protocolo}

A unidade de análise são as 27 UFs no exercício 2024. As
\textit{features} (variáveis preditoras), extraídas do gold dataset, são
três indicadores internos do programa: \textbf{penetração}
(beneficiários\,/\,população\,$\times$\,100), \textbf{per capita}
(\rs/hab em 2021) e \textbf{per beneficiário} (\rs/benef.). Os
\textit{alvos} (variáveis a explicar) são a \textbf{pobreza} (\%) e o
\textbf{IDH-M}. Para o agrupamento, as variáveis são padronizadas em
escore-$z$ (média 0, desvio 1). Para o modelo supervisionado adotamos
validação \textit{leave-one-out} (LOO), apropriada ao $N$ pequeno.
Ferramentas: \texttt{statsmodels} e \texttt{scikit-learn} (Pedregosa et al., 2011).

\subsection{Nível 1 --- correlação e regressão log--log}

A correlação de Pearson mede a associação \textit{linear} entre duas
variáveis, de $-1$ (inversa perfeita) a $+1$ (direta perfeita). Entre
penetração e pobreza obtém-se $r = +0{,}971$; entre per capita e
pobreza, $r = +0{,}971$; entre penetração e IDH-M, $r = -0{,}881$. Em
palavras: \textbf{onde o programa cobre mais gente, há mais pobreza e
menos desenvolvimento}, e a associação é quase perfeita
(Figura \ref{fig:mlscatter}).

\begin{figure}[htbp]
\centering
\includegraphics[width=0.92\linewidth]{fig20-ml-scatter}
\caption{Penetração do Bolsa Família \textit{vs} taxa de pobreza monetária por UF (2024), com cor por \textit{cluster} e reta de ajuste}
\label{fig:mlscatter}
\par\noindent{\small Fonte: elaboração própria, \texttt{ml\_pbf.py}.}
\end{figure}

Para quantificar a elasticidade, estimamos por mínimos quadrados a
regressão $\log(\mathit{pobreza}) = \beta_0 + \beta_1
\log(\mathit{penetra\c{c}\~ao})$:
$$\log(\mathit{pobreza}) = 0{,}714 + 1{,}051 \cdot \log(\mathit{penetra\c{c}\~ao}).$$
Cada número significa o seguinte. O coeficiente $\hat\beta_1 = +1{,}051$
é uma \textbf{elasticidade}: um aumento de 1\,\% na penetração
associa-se a um aumento de cerca de 1,05\,\% na taxa de pobreza
(relação praticamente um-para-um). O \textbf{valor-$p$} $\approx
3{,}8\times10^{-17}$ é a probabilidade de observar associação tão forte
se na verdade não houvesse relação alguma: por ser virtualmente zero, a
associação é estatisticamente inequívoca. O \textbf{R$^2 = 0{,}944$}
indica que 94,4\,\% da variação (em log) da pobreza entre UFs é
acompanhada pela variação da penetração --- isto é, conhecendo
\textit{apenas} a cobertura do programa recupera-se quase toda a
ordenação de pobreza dos estados.

A mesma lógica vale para o desenvolvimento: a regressão $\mathit{IDH\!-\!M}
= 0{,}811 - 0{,}00837\cdot\mathit{penetra\c{c}\~ao}$ ($R^2 = 0{,}776$)
tem coeficiente \textbf{negativo} --- cada ponto percentual adicional de
penetração associa-se a uma queda de $0{,}0084$ no IDH-M ---, confirmando
que o programa se concentra nas UFs menos desenvolvidas.

\begin{table}[H]
\centering
\caption{Modelos de aprendizado de máquina --- o BF como proxy de pobreza e IDH-M (27 UFs, 2024)}
\begin{tabular}{lllr}
\toprule
\textbf{Modelo} & \textbf{Alvo} & \textbf{Métrica} & \textbf{Valor} \\
\midrule
OLS log--log              & pobreza & $\beta_1$ (elasticidade) & $+1{,}051$ \\
                          &         & $R^2$                    & $0{,}944$ \\
OLS                       & IDH-M   & $\beta_1$                & $-0{,}0084$ \\
                          &         & $R^2$                    & $0{,}776$ \\
Regressão linear (LOO)    & pobreza & $R^2$ \,/\, MAE          & $0{,}934$ \,/\, 2{,}61 p.p. \\
Random Forest (LOO)       & pobreza & $R^2$ \,/\, MAE          & $0{,}916$ \,/\, 3{,}17 p.p. \\
\bottomrule
\end{tabular}
\par\noindent{\small Fonte: \texttt{ml\_pbf.py} (\textit{scikit-learn}, \textit{statsmodels}). LOO = validação \textit{leave-one-out}; MAE = erro absoluto médio em pontos percentuais.}
\label{tab:ml}
\end{table}

\subsection{Nível 2 --- agrupamento não supervisionado (K-means)}

O K-means agrupa as UFs por \textit{semelhança} nas quatro variáveis
padronizadas (penetração, per capita, IDH-M, pobreza), sem usar qualquer
rótulo \textit{a priori}. O número de grupos é escolhido pelo
\textit{silhouette}, que varia de $-1$ a $+1$ e mede quão bem cada UF se
encaixa em seu grupo: o máximo ($0{,}695$) ocorre em $k = 2$.

Emergem \textbf{duas Brasis}. O \textbf{Cluster 0} (13 UFs ---
todo o Sul, Sudeste e Centro-Oeste, mais RO e TO) tem pobreza média de
17,2\,\%, IDH-M $0{,}745$, penetração $7{,}8$\,\% e per capita
\rs 460. O \textbf{Cluster 1} (14 UFs --- \textit{todo} o Nordeste e o
Norte) tem pobreza média de 41,8\,\%, IDH-M $0{,}667$, penetração
$17{,}3$\,\% e per capita \rs 1.088. Nenhuma informação de pobreza ou
IDH-M rotulou os grupos de antemão; ainda assim, a partição reproduz
quase exatamente a clivagem socioeconômica conhecida do país
(Figura \ref{fig:mlclusters}).

\begin{figure}[htbp]
\centering
\includegraphics[width=0.92\linewidth]{fig18-ml-clusters}
\caption{Agrupamento K-means das 27 UFs no plano das duas primeiras componentes principais (penetração, per capita, IDH-M e pobreza padronizados)}
\label{fig:mlclusters}
\par\noindent{\small Fonte: elaboração própria, \texttt{ml\_pbf.py}.}
\end{figure}

\subsection{Nível 3 --- modelo supervisionado e validação}

Por fim, testamos se as três variáveis \textit{internas} do programa,
sozinhas, preveem a pobreza. Treinamos dois modelos --- regressão linear
múltipla e \textit{Random Forest} (Breiman, 2001; Hastie et al., 2009)
--- sob validação \textit{leave-one-out}: para cada UF, o modelo aprende
nas outras 26 e prevê a UF retida, de forma que \textbf{nenhum estado é
avaliado com dados que o modelo já viu}. Esse protocolo penaliza
sobreajuste e é o teste honesto de capacidade preditiva.

A regressão linear atinge \textbf{$R^2$(LOO) $= 0{,}934$} e MAE $=2{,}61$
pontos percentuais; a floresta aleatória, $R^2$(LOO) $= 0{,}916$ e MAE
$=3{,}17$. Lendo os números: o $R^2$(LOO) de $0{,}93$ diz que, mesmo
prevendo cada estado \textit{fora da amostra}, as variáveis do BF
explicam 93\,\% da variação da pobreza entre UFs; o MAE de $2{,}6$ p.p.\
é o desvio típico entre a pobreza prevista e a observada
(Figura \ref{fig:mlproxy}). A floresta \textit{não} supera o modelo
linear --- sinal de que a relação é essencialmente linear e que, com
$N = 27$, não há ganho em adicionar flexibilidade. A importância das
variáveis na floresta concentra-se em penetração ($0{,}50$) e per capita
($0{,}47$); o valor por beneficiário ($0{,}03$) é quase irrelevante,
coerente com a homogeneidade do valor médio pago entre UFs após o NBF.

\begin{figure}[htbp]
\centering
\includegraphics[width=0.92\linewidth]{fig19-ml-proxy-fit}
\caption{Pobreza observada (PNAD-C 2019) \textit{vs} prevista pelas variáveis do Bolsa Família (Random Forest, validação \textit{leave-one-out})}
\label{fig:mlproxy}
\par\noindent{\small Fonte: elaboração própria, \texttt{ml\_pbf.py}.}
\end{figure}

\subsection{Síntese e limites}

Os três níveis convergem: a cobertura do Bolsa Família é um \textit{proxy}
empiricamente forte de pobreza ($R^2 \approx 0{,}93$ fora da amostra) e,
inversamente, de desenvolvimento. Isso \textbf{valida quantitativamente}
o uso do programa como indicador territorial na
literatura (Rasella et al., 2013) e, ao mesmo tempo, impõe cautela: trata-se
de associação \textit{transversal} com $N = 27$ e com endogeneidade por
desenho (o programa é \textit{construído} para focalizar a pobreza). As
relações aqui medidas são descritivas e preditivas, \textbf{não causais}.


\section{Discussão}\label{sec:discussao}


\subsection{O salto de 2022\textendash 2023: o que se identifica e o que não}

A Seção \ref{sec:resultados} estabeleceu, descritivamente, dois
saltos de magnitude $> 2\times$ no valor real entre 2021$\to$2022 e
2022$\to$2023. A Seção \ref{sec:causal} mostrou que esses saltos
\textit{não} podem ser atribuídos causalmente à intensidade do
\textit{déficit de cobertura} pré-choque sob a métrica robusta
apropriada para N\,=\,27 \textit{clusters}. O que sobrevive é a
descrição: o programa cresceu uniformemente em todas as UFs,
inclusive nas UFs de menor déficit pré-choque, sugerindo que a
expansão foi \textit{flat} em UF e não um redirecionamento para os
estados com maior gap. Isso é coerente com a literatura
política-administrativa do programa: a expansão veio por elevação
do piso e por inclusão de unipessoais antes restritas, não por
revisão da geografia da elegibilidade.

\subsection{Progressividade decrescente entre regimes}

A descoberta da \textbf{queda da progressividade entre os três
regimes} (K decrescendo monotonicamente em magnitude absoluta de
$|0,334|$ em 2018 para $|0,258|$ em 2024) é distinta dos achados
descritivos da literatura prévia. Sob ônus de Kakwani, isso
significa que a expansão da cobertura entre 2021 e 2024 incorporou
famílias proporcionalmente \textit{menos} concentradas nas UFs de
menor IDH-M. O Auxílio Brasil e o NBF cresceram em todos os
estados, mas cresceram \textit{mais que proporcionalmente} em
estados de IDH-M mediano, atenuando o gradiente progressivo.

\subsection{Sub-cobertura dos estados aparentemente líderes}

O índice de necessidade documenta paradoxo distributivo: UFs no
\textit{topo} do ranking per capita (Maranhão, Acre, Pará, Roraima,
Rondônia) estão de fato \textit{sub-cobertas} em relação à
intensidade da pobreza local ($R < 1$). Em valor por habitante são
campeãs, em \textit{share} relativa à necessidade são deficitárias.
A análise per capita simples, embora correta, esconde essa
informação \textemdash{} é precisamente o tipo de paradoxo que o
índice de necessidade foi desenhado para revelar.

\subsection{Sustentabilidade fiscal sob comparação internacional}

O NBF a 1,10\,\% do PIB é o maior CCT focalizado por renda da
América Latina em magnitude relativa. A Argentina opera AUH a
0,55\,\%, Colômbia 0,40\,\%, México (em sua fase final) 0,30\,\% do
PIB. Esse patamar é replicável apenas pela Bolívia, mas o Renta
Dignidad é universal-por-idade, não focalizado. A trajetória
brasileira de \rs 36,75\,bi (2018) $\to$ \rs 140,88\,bi (2024) em
valores reais coloca o programa em rota de pressão sobre o teto de
gastos sob qualquer cenário de baixo crescimento do PIB. A
discussão pública sobre essa pressão tem ocorrido sem suporte
analítico público em microdados consolidados \textemdash{} este
trabalho contribui esse suporte.

\subsection{PBF vs Emendas: lógicas, formalmente distintas}

A comparação Kakwani (PBF $K \in [-0,33; -0,26]$ ao longo do
período; emendas $K \in [-0,17; -0,13]$) substitui o contraste
descritivo de CV da v1.0. Ambas distribuições são \textit{tecnicamente}
progressivas em relação ao IDH-M, mas o PBF é
sistematicamente mais progressivo \textemdash{} 1,7 a 2,5 vezes a
intensidade. A interpretação é que o PBF responde a critérios
socioeconômicos formais (Cadastro Único, linha de pobreza),
enquanto as emendas respondem à força política das bancadas
estaduais. Ambas as lógicas podem ser legítimas em federação
democrática; o que importa é \textit{tornar explícito} qual
critério rege qual instrumento e \textit{medir} suas implicações
distributivas com aparato cardinal.


\section{Considerações Finais}\label{sec:conclusao}


\subsection{O que documentamos, o que identificamos, o que não}

\textbf{(i) Documentação.} O programa atravessou três regimes
(2003\textendash 2021, 2022, 2023\textendash 2025). Em valores
reais, quase quadruplicou entre 2018 e 2024 (\rs 36,75\,bi $\to$
\rs 140,88\,bi, $+283,4\,\%$). O número de beneficiários cresceu
$36,4\,\%$ (16,99\,M $\to$ 23,18\,M, com pico em 2023 de 24,51\,M).
O valor por beneficiário cresceu $180,9\,\%$ ($\rs 2.163 \to
\rs 6.078$). A focalização territorial em 2025 apresenta razão
$6,98:1$ entre os extremos, com Nordeste a \rs 1.051/hab e Sul a
\rs 283/hab.

\textbf{(ii) Identificação.} O DiD~2$\times$2 e o TWFE
clusterizado sobre MP 1.061/2021 produzem coeficientes grandes e
estatisticamente significativos sob HC3 ($\hat\beta \approx +220$
\rs/hab, $p < 0,001$); o \textit{wild-cluster bootstrap},
contudo, devolve $p_{\mathrm{WCB}} = 0,49$, e o teste de
\textit{parallel trends} pré-tratamento é rejeitado. \textbf{O
salto descritivo é incontroverso; a atribuição \textit{causal}
desse salto à intensidade do déficit de cobertura pré-choque não é
sustentada pela métrica robusta apropriada para N=27}. Esse é o
ônus do desenho quasi-experimental sobre painel UF\,$\times$\,Ano
com poucos \textit{clusters} e tendências pré-divergentes
\textemdash{} resultado honesto que reflete os limites da
identificação possível com os dados publicamente disponíveis.

\textbf{(iii) Equidade.} PBF/NBF é progressivo
($K = -0,26$ em 2024, IC 95\,\% bootstrap $[-0,30; -0,19]$) e
sistematicamente mais progressivo que emendas
($K = -0,15$ em 2024). A progressividade do PBF, contudo, decresceu
monotonicamente entre 2018 e 2024 \textemdash{} a expansão
incorporou famílias proporcionalmente menos concentradas nas UFs
de menor IDH-M.

\textbf{(iv) Necessidade.} UFs aparentemente \textit{campeãs} em
per capita (MA, RR, RO, MA) estão sub-cobertas em relação à
intensidade da pobreza local ($R < 1$). UFs ricas têm
sobre-cobertura relativa.

\textbf{(v) Benchmark internacional.} O NBF é o CCT de maior valor
per beneficiário e maior cobertura da América Latina, com custo
relativo (\%PIB) elevado para o padrão regional.

\subsection{Agenda de pesquisa}

(a) Ampliação para microdados \textit{municipais} (não apenas UF)
para mitigar o problema de poucos \textit{clusters} no DiD;
(b) integração com dados administrativos do Cadastro Único para
RDD estrutural sobre a linha de pobreza
($y_{\mathrm{Cad\acute{U}n}} \lessgtr \rs 218$);
(c) cruzamento com indicadores de pobreza multidimensional (IPM,
IDH-M atualizado pelo PNUD); (d) comparação prospectiva com
programas congêneres em outros países da América Latina sob
métricas harmonizadas (CEPAL e World Bank ASPIRE).

\subsection{Contribuição metodológica}

A principal contribuição metodológica é um \textit{pipeline} de dados
\textbf{aberto e gratuito}, construído inteiramente sobre o
\textbf{Databricks Free Edition} --- a camada sem custo da plataforma ---,
o que permite que qualquer pessoa reproduza a consolidação dos mais de
2,2 bilhões de registros \textbf{sem investimento em infraestrutura}.
A arquitetura medallion (bronze \textrightarrow{} silver
\textrightarrow{} gold) parte dos arquivos mensais brutos da CGU e
entrega um \textit{gold dataset} público em painel UF\,$\times$\,Ano,
acompanhado de uma suíte de testes automatizados de qualidade de dados
(invariantes de esquema, contagem, plausibilidade, monotonia
populacional e focalização Nordeste\,$>$\,Sul) e aderente aos
princípios FAIR (Wilkinson et al., 2016). Todos os indicadores
descritivos, os índices de equidade, os modelos de aprendizado de
máquina e as figuras deste artigo são gerados a partir desse
\textit{gold dataset} por \textit{scripts} reproduzíveis, sem nenhum
valor \textit{hardcoded}, de modo que terceiros podem auditar e
estender a análise a custo marginal praticamente nulo.


\section*{Referências}
\begin{singlespace}\small
\newcommand{\refitem}{\par\smallskip\noindent\hangindent=1.25cm\hangafter=1\relax}

\refitem Armbrust, M., Ghodsi, A., Xin, R., \& Zaharia, M. (2021). Lakehouse: A new generation of open platforms that unify data warehousing and advanced analytics. In \textit{11th Annual Conference on Innovative Data Systems Research (CIDR)}.

\refitem Banco Central do Brasil. (2024). \textit{Sistema Gerenciador de Séries Temporais --- série 433 (IPCA)} [Conjunto de dados]. \url{https://api.bcb.gov.br/dados/serie/bcdata.sgs.433/dados}

\refitem Breiman, L. (2001). Random forests. \textit{Machine Learning, 45}(1), 5--32.

\refitem Cameron, A. C., Gelbach, J. B., \& Miller, D. L. (2008). Bootstrap-based improvements for inference with clustered errors. \textit{The Review of Economics and Statistics, 90}(3), 414--427.

\refitem Campello, T., \& Neri, M. C. (Orgs.). (2013). \textit{Programa Bolsa Família: uma década de inclusão e cidadania}. IPEA.

\refitem CEPAL --- Comisión Económica para América Latina y el Caribe. (2024). \textit{Panorama social de América Latina y el Caribe 2024}. CEPAL. \url{https://www.cepal.org/es/publicaciones/panoramasocial}

\refitem CGU --- Controladoria-Geral da União. (2025). \textit{Portal da Transparência: microdados de pagamentos do Bolsa Família, Auxílio Brasil e Novo Bolsa Família} [Conjunto de dados]. \url{https://portaldatransparencia.gov.br/download-de-dados/bolsa-familia-pagamentos}

\refitem Fiszbein, A., \& Schady, N. (2009). \textit{Conditional cash transfers: Reducing present and future poverty}. World Bank.

\refitem Hastie, T., Tibshirani, R., \& Friedman, J. (2009). \textit{The elements of statistical learning} (2nd ed.). Springer.

\refitem Hoffmann, R. (2006). Transferências de renda e a redução da desigualdade no Brasil e cinco regiões entre 1997 e 2004. \textit{Econômica, 8}(1), 55--81.

\refitem IBGE --- Instituto Brasileiro de Geografia e Estatística. (2019). \textit{Pesquisa Nacional por Amostra de Domicílios Contínua (PNAD-C) --- Tabela 6688} [Conjunto de dados]. \url{https://sidra.ibge.gov.br/tabela/6688}

\refitem IBGE --- Instituto Brasileiro de Geografia e Estatística. (2024). \textit{SIDRA --- Tabela 6579: população residente estimada} [Conjunto de dados]. \url{https://sidra.ibge.gov.br/tabela/6579}

\refitem Kakwani, N. C. (1977). Measurement of tax progressivity: An international comparison. \textit{The Economic Journal, 87}(345), 71--80.

\refitem Kukreja, M. (2021). \textit{Data engineering with Apache Spark, Delta Lake, and Lakehouse}. Packt Publishing.

\refitem Paes-Sousa, R., Santos, L. M. P., \& Quaresma, J. M. (Orgs.). (2011). \textit{Bolsa Família: avaliação de impacto e o desafio da intersetorialidade}. Ministério do Desenvolvimento Social.

\refitem Pedregosa, F., Varoquaux, G., Gramfort, A., Michel, V., Thirion, B., Grisel, O., \ldots Duchesnay, É. (2011). Scikit-learn: Machine learning in Python. \textit{Journal of Machine Learning Research, 12}, 2825--2830.

\refitem PNUD, IPEA, \& Fundação João Pinheiro. (2013). \textit{Atlas do desenvolvimento humano no Brasil}. \url{http://atlasbrasil.org.br/}

\refitem Rasella, D., Aquino, R., Santos, C. A. T., Paes-Sousa, R., \& Barreto, M. L. (2013). Effect of a conditional cash transfer programme on childhood mortality: A nationwide analysis of Brazilian municipalities. \textit{The Lancet, 382}(9886), 57--64.

\refitem Rocha, S. (2008). Transferências de renda federais: focalização e impactos sobre pobreza e desigualdade. \textit{Revista de Economia Contemporânea, 12}(1), 67--96.

\refitem Roth, J. (2022). Pretest with caution: Event-study estimates after testing for parallel trends. \textit{American Economic Review: Insights, 4}(3), 305--322.

\refitem Soares, F. V. (2010). \textit{Bolsa Família, its design, its impacts and possibilities for the future} (Working Paper No. 137). International Policy Centre for Inclusive Growth.

\refitem Wilkinson, M. D., Dumontier, M., Aalbersberg, I. J., Appleton, G., Axton, M., Baak, A., \ldots Mons, B. (2016). The FAIR Guiding Principles for scientific data management and stewardship. \textit{Scientific Data, 3}, 160018. \url{https://doi.org/10.1038/sdata.2016.18}

\refitem World Bank. (2024a). \textit{ASPIRE: Atlas of Social Protection Indicators of Resilience and Equity} [Conjunto de dados]. \url{https://www.worldbank.org/en/data/datatopics/aspire}

\refitem World Bank. (2024b). \textit{International Comparison Program (ICP) 2021} [Conjunto de dados]. \url{https://www.worldbank.org/en/programs/icp}

\refitem Zaharia, M., Xin, R., Wendell, P., Das, T., Armbrust, M., Dave, A., \ldots Stoica, I. (2016). Apache Spark: A unified engine for big data processing. \textit{Communications of the ACM, 59}(11), 56--65.
\end{singlespace}


\end{document}
